% Requires amsmath, booktabs, multirow, and graphicx.
\section{Additional details for the electronic AKI alert trial}
\label{supp:trial3}

\subsection{Endpoint construction and variable mapping}

The application to the trial of electronic health record alerts for AKI
\citep{wilson2021electronic} used the common restriction horizon $S=14$ days after
randomization. The treatment and endpoint mappings were verified against the accompanying
analysis data dictionary before analysis:
\[
\begin{array}{lll}
\text{Role} & \text{Dataset variable} & \text{Coding}\\
\hline
\text{Treatment} & \texttt{Treatment}
  & A_i=1:\ \texttt{Alert};\quad A_i=0:\ \texttt{Usual Care}\\
\text{Death by day 14} & \texttt{YS\_death\_event} & 1=\text{event}\\
\text{Dialysis by day 14} & \texttt{YS\_dialysis\_event} & 1=\text{event}\\
\text{AKI progression by day 14} & \texttt{YS\_aki\_progression\_event}
  & 1=\text{event}.
\end{array}
\]
The three component variables were complete binary 0/1 indicators, and their elementwise
maximum exactly reproduced the trial-defined 14-day composite
\texttt{YP\_composite\_event}. Although the dataset also contained
\texttt{YP\_composite\_time}, it was not used because the application concerned the
fixed-$S$ status record
\[
  \mathcal D_i(14)
  =
  \{\text{death by day 14},\text{dialysis by day 14},
    \text{AKI progression by day 14}\}.
\]

For a treatment--control pair $(i,j)$, define
\[
  E_{q-1}(i,j)
  =\prod_{r=1}^{q-1}I\{D_{ri}(14)=D_{rj}(14)\},
  \qquad E_0(i,j)=1.
\]
The win, loss, and tie kernels used in the analysis were
\[
\begin{split}
  \varphi_w\{\mathcal D_i(14),\mathcal D_j(14)\}
    &=\sum_{q=1}^3 E_{q-1}(i,j)
      I\{D_{qi}(14)<D_{qj}(14)\},\\
  \varphi_l\{\mathcal D_i(14),\mathcal D_j(14)\}
    &=\sum_{q=1}^3 E_{q-1}(i,j)
      I\{D_{qi}(14)>D_{qj}(14)\},\\
  \varphi_t\{\mathcal D_i(14),\mathcal D_j(14)\}
    &=\prod_{q=1}^3 I\{D_{qi}(14)=D_{qj}(14)\}.
\end{split}
\]
Because only the first unequal component contributes, these kernels implement the
prespecified priority death $\succ$ dialysis $\succ$ AKI progression.

\subsection{Baseline adjustment variables}

No participant was removed because of a missing baseline covariate, and no missing value
was imputed. Table~\ref{tab:trial3_covariates} records the prespecified eligibility
decisions. Hospital was treated as a six-level factor, with one reference category and five
dummy variables. All binary variables were reference coded. The centered LADJ design
therefore contained 11 columns.

\begin{table}[t]
\centering
\small
\caption{Prespecified baseline covariates for the AKI alert trial LADJ analysis.}
\label{tab:trial3_covariates}
\begin{tabular}{lllrl}
\toprule
Covariate & Dataset variable & Type & Missing & Primary LADJ \\
\midrule
Hospital/site & \texttt{X\_hospital\_0m} & Categorical & 0 & Included \\
Age & \texttt{X\_age\_0m} & Continuous & 507 & Excluded \\
Age over 90 & \texttt{X\_age\_over\_90\_0m} & Binary & 0 & Included \\
Sex & \texttt{X\_sex\_0m} & Binary & 0 & Included \\
ICU status & \texttt{X\_icu\_0m} & Binary & 0 & Included \\
Serum creatinine & \texttt{X\_creat\_0m} & Continuous & 2 & Excluded \\
Baseline AKI stage & Not available & -- & -- & Excluded \\
Elixhauser score & \texttt{X\_elx\_score\_0m} & Continuous & 0 & Included \\
Diabetes & \texttt{X\_diabetes\_0m} & Binary & 0 & Included \\
Congestive heart failure & \texttt{X\_chf\_0m} & Binary & 0 & Included \\
\bottomrule
\end{tabular}
\end{table}

\subsection{Descriptive endpoint results}

Table~\ref{tab:trial3_events} reports both the component event rates and mutually exclusive
severity categories. For the latter, participants with overlapping events were assigned to
their highest-priority observed event.

\begin{table*}[t]
\centering
\small
\caption{Fourteen-day endpoint summaries by randomized arm in the AKI alert trial. Values are
$n$ (\%).}
\label{tab:trial3_events}
\begin{tabular}{lrrrrrrrr}
\toprule
Arm & $N$ & Death & Dialysis & AKI progression &
Death category & Alive+dialysis & Alive+AKI & None \\
\midrule
Usual Care & 2971 & 265 (8.92) & 93 (3.13) & 461 (15.52) &
265 (8.92) & 65 (2.19) & 292 (9.83) & 2349 (79.06) \\
Alert & 3059 & 272 (8.89) & 106 (3.47) & 487 (15.92) &
272 (8.89) & 60 (1.96) & 321 (10.49) & 2406 (78.65) \\
\bottomrule
\end{tabular}
\end{table*}

Among $9{,}088{,}289$ treatment--control pairs, $16.23\%$ were decided at the death
tier, $3.88\%$ at the dialysis tier, and $16.27\%$ at the AKI-progression tier; the
remaining $63.62\%$ were ties. The resulting unadjusted elementary estimates were
\[
  \bar\tau_w=0.180156,\qquad
  \bar\tau_l=0.183652,\qquad
  \bar\tau_t=0.636192.
\]
After linear adjustment,
\[
  \widehat\tau_w^{\mathrm{lin}}=0.178194,\qquad
  \widehat\tau_l^{\mathrm{lin}}=0.186105,\qquad
  \widehat\tau_t^{\mathrm{lin}}=0.635701.
\]

\subsection{Complete treatment-effect results}

Table~\ref{tab:trial3_full_results} reports inference on the NB, LWR, and LWO scales and
the corresponding exponentiated WR and WO scales. Wald tests for WR and WO were conducted
on the LWR and LWO scales, respectively. Relative efficiency was calculated as
$\widehat{\mathrm{ASE}}_{\mathrm{UNADJ}}^2/
 \widehat{\mathrm{ASE}}_{\mathrm{LADJ}}^2$ on the analysis scale.

\begin{table*}[t]
\centering
\scriptsize
\caption{Complete unadjusted and LADJ results for the AKI alert trial.}
\label{tab:trial3_full_results}
\setlength{\tabcolsep}{4pt}
\resizebox{\textwidth}{!}{%
\begin{tabular}{llrrrrrrr}
\toprule
Measure & Method & Estimate & ASE & 95\% CI & Natural-scale estimate (95\% CI) &
$p$-value & Relative efficiency & PVR, \% \\
\midrule
\multirow{2}{*}{NB}
& $\bar\eta_m$ & -0.00350 & 0.01060 & $(-0.02428,\ 0.01728)$ & -- &
0.7416 & 1.000 & -- \\
& $\widehat\eta_m^{\mathrm{lin}}$ & -0.00791 & 0.01010 &
$(-0.02772,\ 0.01190)$ & -- & 0.4337 & 1.101 & 9.16 \\
\addlinespace[2pt]
\multirow{2}{*}{LWR / WR}
& $\bar\eta_m$ & -0.01922 & 0.05828 & $(-0.13346,\ 0.09501)$ &
$0.98096\ (0.87507,\ 1.09967)$ & 0.7415 & 1.000 & -- \\
& $\widehat\eta_m^{\mathrm{lin}}$ & -0.04344 & 0.05549 &
$(-0.15220,\ 0.06533)$ & $0.95749\ (0.85882,\ 1.06751)$ &
0.4338 & 1.103 & 9.35 \\
\addlinespace[2pt]
\multirow{2}{*}{LWO / WO}
& $\bar\eta_m$ & -0.00699 & 0.02120 & $(-0.04855,\ 0.03457)$ &
$0.99303\ (0.95261,\ 1.03517)$ & 0.7416 & 1.000 & -- \\
& $\widehat\eta_m^{\mathrm{lin}}$ & -0.01582 & 0.02021 &
$(-0.05544,\ 0.02379)$ & $0.98430\ (0.94607,\ 1.02408)$ &
0.4338 & 1.101 & 9.15 \\
\bottomrule
\end{tabular}%
}
\end{table*}

The LADJ Gram matrix had full rank (11 of 11). Its condition number was 271.4 on the raw
feature scale and 4.20 after feature scaling. The closed-form imbalance correction and
direct projection formula agreed to numerical precision, and all adjusted estimates were
exact transformations of the adjusted elementary win--loss vector.
